\Needspace{14\baselineskip}
\section{Related Work}
\label{sec:related-work}

% PRISMA Item 3: Rationale — positioning against prior reviews

This section positions the present review against two classes of
secondary work: surveys on memory for LLM-based agents in general and
systematic reviews on LLMs applied to software engineering.
The first class is represented mainly by the recent survey of
\citet{du2026memory};
the second groups five SLRs published between 2024 and 2025.
Du covers memory for LLM agents broadly, including within-session
context management, rather than cross-session persistence as a
distinct object;
the five SE-focused SLRs treat memory only partially or not at all.
Neither class adopts cross-session persistence as the primary object
while also restricting scope to coding agents, and none reports
per-task operational cost analysis.
\autoref{tab:slrs-previas} summarises the comparison by thematic focus
and coverage of these two dimensions.

\begin{table}[htbp]
\centering
\caption{Comparison between related secondary work and this review.
\emph{Persist.} marks whether the work treats cross-session knowledge
persistence as a primary object of analysis (``Yes''), as one of several
components (``Partial''), or not at all (``No'').
\emph{Cost} marks whether the work reports per-task operational cost
data for the systems it reviews.}
\label{tab:slrs-previas}
\small
\setlength{\tabcolsep}{4pt}
\begin{tabularx}{\textwidth}{@{} l c >{\raggedright\arraybackslash}X c c @{}}
\toprule
\textbf{Work} & \textbf{Year} & \textbf{Main focus} &
  \textbf{Persist.} & \textbf{Cost} \\
\midrule
\citet{du2026memory} & 2026 &
  Memory for LLM agents in general &
  Partial & No \\
\citet{hou2024llmse} & 2024 &
  LLMs for software engineering (395 studies) &
  No & No \\
\citet{zhang2024unifying} & 2024 &
  Language models for code (NLP $\cap$ SE) &
  No & No \\
\citet{jin2024llmagents} & 2024 &
  Evolution of LLMs into agents in SE &
  No & No \\
\citet{wang2024llmagents} & 2024 &
% audit:ignore-next-line — describes cited work's own framing
  Autonomous agents for SE &
  Partial & No \\
\citet{liu2024llmagents} & 2024 &
  LLM agents for software engineering (117 studies) &
  Partial & No \\
\midrule
\textbf{This review} & 2026 &
  \textbf{Knowledge persistence in AI-based coding agents
  (28 studies)} &
  \textbf{Yes} & \textbf{Yes} \\
\bottomrule
\end{tabularx}
\end{table}

\subsection{Surveys on memory for LLM agents}
\label{sec:rel-memoria}

\citet{du2026memory} presents a recent survey on memory
mechanisms for LLM-based agents, proposing a three-dimensional taxonomy
(temporal scope, representational substrate, and control policy) that
this review adopts and refines as the starting point of SQ1.
The survey catalogues five families of memory mechanisms: resident
context compression, retrieval-augmented stores, self-correcting
reflective memory, virtual hierarchical contexts, and policy-learned
memory management.
The survey's scope is explicitly broad.
It covers general agents for web navigation, mathematical reasoning and
dialogue, without restriction to the coding domain.
The coverage of software engineering in Du's survey is therefore
limited to ChatDev and MetaGPT.
Du treats memory as a primary topic, but not cross-session persistence
in coding agents as a bounded review object.
Systems central to the present investigation are not mentioned:
MemCoder~\citep{deng2026memcoder},
CodeMEM~\citep{wang2026codemem}, and the Repository Memory approach of
\citet{wang2026improving}.
The present review complements that work by restricting the scope to
coding agents and by incorporating per-task operational cost analysis,
a dimension absent from Du's survey.

\subsection{Systematic reviews on LLMs in software engineering}
\label{sec:rel-llmse}

Five recent systematic reviews map LLM applications in software
engineering with broad scope.
\citet{hou2024llmse} cover 395 studies published up to 2023, organised
by engineering task (code generation, repair, testing, maintenance) and
by model type, with no subsection dedicated to memory or cross-session
persistence.
\citet{zhang2024unifying} unify natural language processing and
software engineering perspectives on models for code, with emphasis on
pre-training and evaluation benchmarks.
\citet{jin2024llmagents} trace the transition from LLMs to LLM-based
agents in software engineering but treat memory as a generic
architectural component, without characterising cross-session
persistence architectures.
\citet{wang2024llmagents} survey autonomous agents in software
engineering with a focus on execution pipelines, and
\citet{liu2024llmagents} cover 117 studies with partial discussion of
memory modules, but without cross-system synthesis of persistence
effectiveness and cost.
None of these reviews directly answers SQ1--SQ5 of this investigation.
The missing dimensions are the taxonomy of persistence architectures,
the comparability of memory-specific benchmarks, the quantification of
gains over no-memory baselines, the cost-effectiveness relationship,
and the mechanisms by which memory can degrade performance.

\subsection{Positional gap of this review}
\label{sec:rel-lacuna}

This review restricts the scope to AI-based coding agents and adopts
cross-session knowledge persistence as the primary object of
investigation, not as an ancillary architectural component of a broader
synthesis.
On this basis, it incorporates per-task operational cost analysis as a
dedicated dimension (SQ4).
Cost is where this review departs most sharply from prior secondary
work.
None of the previously identified SLRs simultaneously gathers a focus
on cross-session persistence, a coding-agent scope, and cross-study cost
quantification.
This scope isolates two practical questions:
what does memory improve, and what does it cost?
It also exposes the absence of head-to-head comparisons between
architectures and the mechanisms by which memory can harm agent
performance.
